\section*{第 \ref{ch:b1:membranes-transport} 章——膜与膜运输}

\begin{solution}{exo:b1:membranes-transport:1}
由磷脂组成的双层，尾部朝内，并含胆固醇，形成一层
\qty{5}{nm} 厚的薄片。蛋白质横跨其中或附着于一侧，
各自独立。脂质和大多数蛋白质可在平面内侧向扩散，
平面内，属二维流体。两层小叶的脂质组成不同，
糖链仅位于外表面。
\end{solution}

\begin{solution}{exo:b1:membranes-transport:2}
$\mathrm{O_2}$（小，非极性）$>$ 乙醇（小，弱极性）
$>$ 水（小，极性）$\gg$ 葡萄糖（大，极性）$\gg$
$\mathrm{Na^+}$（带电，水合）。通透性随极性、
大小，尤其是电荷而降低，因为溶质必须溶入
疏水核心。
\end{solution}

\begin{solution}{exo:b1:membranes-transport:3}
蔗糖：$-2.48\times 0.1 = \qty{-0.25}{MPa}$。NaCl 解离为
两个离子，即 \qty{0.2}{osmol/L}：\qty{-0.50}{MPa}。
\end{solution}

\begin{solution}{exo:b1:membranes-transport:4}
$E_{Cl} = (61.5/(-1))\log(120/4) = -61.5\times 1.48 = \qty{-91}{mV}$.
在 \qty{-89}{mV} 时，氯离子与\omterm{thm:b1:membranes-transport:nernst}{平衡电位}相差不超过 \qty{2}{mV}：它是
被动分布的。
\end{solution}

\begin{solution}{exo:b1:membranes-transport:5}
\omterm{def:b1:cell-unit-of-life:cell}{细胞}表面积 $\num{4.7e9}\times 10^{-10}\,\mathrm{m^2} =
\qty{0.47}{m^2}$；比值 $0.92/0.47 = 1.96 \approx 2$：脂质覆盖
\omterm{def:b1:cell-unit-of-life:cell}{细胞}两次，因此膜为双分子层。
\end{solution}

\begin{solution}{exo:b1:membranes-transport:6}
\omterm{def:b1:cell-unit-of-life:cell}{细胞} $\Psi = -0.9 + 0.5 = \qty{-0.4}{MPa}$；溶液 \qty{-0.6}{MPa}：
水离开\omterm{def:b1:cell-unit-of-life:cell}{细胞}。平衡时 $\Psi_{\text{细胞}} = -0.6$：
$\Psi_p = -0.6 + 0.9 = \qty{0.3}{MPa}$；\omterm{def:b1:cell-unit-of-life:cell}{细胞}失去部分\omterm{def:b1:eukaryotic-cell:plantcell}{膨压}，但
仍保持膨胀状态。
\end{solution}

\begin{solution}{exo:b1:membranes-transport:7}
浓度由 10 增至 20 时，速率仅提高
\qty{15}{\%}：已达饱和。以 $1/v$ 对 $1/c$ 作图（或注意到
$v = 1.0$ 对应 \qty{5}{mmol/L}，$1.67$ 为 $c$ 趋于无穷时的值），可得
$V_{\max} \approx \qty{1.7}{\micro mol/min}$，半最大速率对应的浓度约为
\qty{3}{mmol/L}：这是载体介导的\omterm{def:b1:membranes-transport:transporters}{协助扩散}。
\end{solution}

\begin{solution}{exo:b1:membranes-transport:8}
静息时，膜的通透性主要由开放的钾离子通道决定；
钾离子外流，直至膜内负电性足够强，足以
阻止其继续外流，即接近 $E_K$。若钠通道开放，通透性
将主要由钠离子决定，电位将向
$E_{Na}$ 偏移，约为 \qty{+60}{mV}——即动作电位。
\end{solution}

\begin{solution}{exo:b1:membranes-transport:9}
每摩尔 $\mathrm{Na^+}$ 内流：$RT\ln(12/145) + F(-0.060) = -6.4 - 5.8
= \qty{-12.2}{kJ}$；两个则为 \qty{-24.4}{kJ}。葡萄糖逆浓度梯度运输的代价为
$RT\ln(c_{\text{内}}/c_{\text{外}})$；令二者相等，得
$\ln(c_{\text{内}}/c_{\text{外}}) = 24\,400/2577 = 9.5$，比值
约为 \num{13000}。
\end{solution}

\begin{solution}{exo:b1:membranes-transport:10}
流动性下降（尾部紧密堆积，双层趋近凝胶）；氧气
扩散减慢但仍在进行；作为酶的\omterm{def:b1:membranes-transport:transporters}{钠钾泵}几乎停止工作；
渗漏仍在继续，数小时后钠离子内流、钾离子外流，
浓度梯度逐渐消失；电位随之趋于零，
梯度消退；钠离子内流，而蛋白质的渗透拉力不再
被平衡，水分进入，\omterm{def:b1:cell-unit-of-life:cell}{细胞}肿胀——低温保存的\omterm{def:b1:cell-unit-of-life:cell}{细胞}
与\omterm{def:b1:mammal-organization:organ}{器官}正是因此而肿胀。
\end{solution}

\begin{solution}{exo:b1:membranes-transport:11}
在 \qty{37}{\celsius} 下混合、在 \qty{4}{\celsius} 下不混合，说明这是一个
依赖温度的过程——即一种在低温下会变成固体的流体中的扩散。
若无 ATP 时也不发生混合，则说明存在主动的、耗能的
重新分布（马达蛋白、囊泡循环），而非
扩散。修正后的结论是混合无需 ATP，这正是
\omterm{def:b1:membranes-transport:membrane}{流动镶嵌}模型所要求的：蛋白质借热扩散在二维
流体中运动，\omterm{def:b1:cell-unit-of-life:cell}{细胞}无需做功。
\end{solution}

\begin{solution}{exo:b1:membranes-transport:12}
电位是电场力与通透性最强的离子的扩散相平衡时的值，
由能斯特给出；\omterm{def:b1:cell-unit-of-life:cell}{细胞}借泵建立梯度，
借通道设定通透性，
电位随之确定。一层 \qty{1}{\micro F/cm^2} 的膜在
\qty{-70}{mV} 下携带 $\qty{7e-8}{C/cm^2}$ 的电荷，约合每
平方厘米 $\num{4e11}$ 个离子——对于 \qty{20}{\micro m} 的\omterm{def:b1:cell-unit-of-life:cell}{细胞}，约 $10^5$ 个离子，
与\omterm{def:b1:cell-unit-of-life:cell}{细胞}内 $10^{10}$ 个钾离子相比，充电所需的离子只占极小比例，
且只移动极短距离，即可为电容器充电。
浓度梯度是储存，电位则是其读数；
泵的直接贡献（排出三个电荷、流入两个）仅为
几毫伏。
\end{solution}

\begin{solution}{pb:b1:membranes-transport:1}
\textbf{1.} $E_{Na} = 26.7\ln(145/12) = \qty{+66}{mV}$; $E_K =
26.7\ln(5/140) = \qty{-89}{mV}$.
\textbf{2.} $\mathrm{K^+}$（距平衡 \qty{29}{mV}）比
$\mathrm{Na^+}$（距平衡 \qty{126}{mV}）更远：膜对
$\mathrm{K^+}$ 的通透性要大得多。
\textbf{3.} $\mathrm{Na^+}$ 内流：$RT\ln(12/145) + F(-0.060) = -6.4 -
5.8 = \qty{-12.2}{kJ/mol}$。$\mathrm{K^+}$ 外流：$RT\ln(5/140) -
F(-0.060) = -8.6 + 5.8 = \qty{-2.8}{kJ/mol}$.
\textbf{4.} 每个循环的代价 $= 3\times 12.2 + 2\times 2.8 = \qty{42.2}{kJ}$，
与每个 ATP 提供的 \qty{50}{kJ} 相比，效率为 \qty{84}{\%}。
\textbf{5.} 泵使胞质钠离子保持低浓度；顶端膜上的\omterm{prop:b1:membranes-transport:secondary}{同向转运体}
利用向内的钠离子梯度，将葡萄糖从肠腔
拉入\omterm{def:b1:cell-unit-of-life:cell}{细胞}；泵必须位于另一侧膜上，使其排出的钠离子
进入血液，而非回到肠腔，否则肠腔中的钠离子将
被循环利用，净运输将为零。
\textbf{6.} 每摩尔葡萄糖需 $2\times 12.2 = \qty{24.4}{kJ/mol}$。
\textbf{7.} $\Delta G = RT\ln(c_C/c_L)$（葡萄糖不带电）。
\textbf{8.} $\ln(c_C/c_L) = 24\,400/2577 = 9.47$；$c_C/c_L =
\num{1.3e4}$.
\textbf{9.} 原则上可达 \qty{1300}{mmol/L}：该比值
永远不会达到，因为葡萄糖经基底外侧的载体离开\omterm{def:b1:cell-unit-of-life:cell}{细胞}；
实际上胞质浓度保持在 \qtyrange{5}{20}{mmol/L} 附近。
\textbf{10.} 胞质（高于 \qty{5}{mmol/L}）中的葡萄糖浓度高于
血液；载体使葡萄糖顺此梯度
被动扩散进入血液。
\textbf{11.} 若方向相反，\omterm{prop:b1:membranes-transport:secondary}{同向转运体}会将葡萄糖从血液
泵入\omterm{def:b1:cell-unit-of-life:cell}{细胞}，而载体则会使其漏入肠腔：
肠道将分泌葡萄糖。转运体的极性
决定了吸收的方向。
\textbf{12.} \omterm{prop:b1:membranes-transport:secondary}{同向转运体}仅在葡萄糖存在时才摄取钠离子；因此，
肠腔中的葡萄糖因此仍能驱动钠离子经此途径吸收，即使
霍乱毒素阻断了其他途径；氯离子在电性上跟随钠离子，
水分则依盐分的\omterm{def:b1:membranes-transport:osmosis}{渗透作用}随之进入。
\textbf{13.} 每秒 $\qty{4e-15}{mol/s}$ 的 $\mathrm{Na^+}$，即
$\num{2.4e9}$ 个离子；按每个循环三个计，
每秒需 $\num{8e8}$ 个循环。
\textbf{14.} 每秒 $\num{8e8}$ 个 ATP $= \qty{1.3e-15}{mol/s}$；按
每个葡萄糖 30 个 ATP 计，则为 \qty{4.4e-17}{mol/s} 的葡萄糖，
占所吸收葡萄糖的 \qty{2.2}{\%}。
\textbf{15.} 每秒 $\num{1.2e9}$ 个葡萄糖分子，按每个\omterm{prop:b1:membranes-transport:secondary}{同向转运体} 50 个计，
需 $\num{2.4e7}$ 个\omterm{prop:b1:membranes-transport:secondary}{同向转运体}；膜面积 \qty{500}{\micro m^2}，
即每平方微米 \num{48000} 个，亦即每 \qty{20}{nm} 见方一个——
一层转运体排列十分密集的膜。
\textbf{16.} 进入的渗透摩尔数：$4\times\qty{2e-15}{mol/s} =
\qty{8e-15}{osmol/s}$；在 \qty{0.3}{osmol/L} 时，水的
流量 $= \qty{2.7e-14}{L/s} = \qty{27}{\micro m^3/s}$。\omterm{def:b1:cell-unit-of-life:cell}{细胞}体积
$25\times 5\times 5 = \qty{625}{\micro m^3}$：\qty{23}{s} 内加倍。
水分离开\omterm{def:b1:cell-unit-of-life:cell}{细胞}的速度与进入的速度相同，经由基底外侧
膜流出，这正是肠道吸收水分的方式。
\textbf{17.} $\qty{2.7e-14}{L/s}$ 的水相当于 \qty{1.5e-12}{mol/s}；与每个葡萄糖（$\qty{2e-15}{mol/s}$）相比，约为 750 个水分子。
\textbf{18.} $S = 2(25\times 5 + 25\times 5 + 5\times 5) =
\qty{550}{\micro m^2} = \qty{5.5e-6}{cm^2}$; $C = \qty{5.5e-12}{F}$.
\textbf{19.} $Q = CV = 5.5\times 10^{-12}\times 0.060 =
\qty{3.3e-13}{C}$；$/\num{1.6e-19} = \num{2.1e6}$ 个离子。
\textbf{20.} $\num{2.1e6}/550 = 3800$ 个离子每平方微米，即每 \qty{16}{nm} 见方的膜上有一个多余电荷。
\textbf{21.} \omterm{def:b1:cell-unit-of-life:cell}{细胞}内的 $\mathrm{K^+}$：$0.140\,\mathrm{mol/L}\times
\qty{6.25e-13}{L}\times\num{6e23} = \num{5.3e10}$ 个离子；充电所需
的电荷仅占其中 \qty{0.004}{\%}：给膜充电
并不改变浓度。
\textbf{22.} 每个循环净输出 $+e$ 的电荷：电流 $\num{8e8}\times
\num{1.6e-19} = \qty{1.3e-10}{A}$；一秒内膜将获得
\qty{1.3e-10}{C} 的电荷，即 $\Delta V = Q/C = \qty{24}{V}$。实际并非如此，
因为每个被泵出的电荷都会立即被通过通道
回流的离子所平衡（主要是 $\mathrm{K^+}$ 内漏减少，
或 $\mathrm{Na^+}$ 重新进入）：膜的电导使
泵的电流短路，泵只造成毫伏级的变化。
\textbf{23.} 胞质钠离子在数分钟内升高；\omterm{prop:b1:membranes-transport:secondary}{同向转运体}失去
驱动力，葡萄糖吸收停止；电位随之衰减，
趋于零，浓度梯度逐渐耗尽；钠离子与水分进入，
\omterm{def:b1:cell-unit-of-life:cell}{细胞}肿胀。
\textbf{24.} 阻断钾离子渗漏后，钠离子渗漏
占主导：电位趋向 $E_{Na}$，负值大为减小
（去极化）。
\textbf{25.} 胞质与肠腔之比约为 \num{13000}，由
每个葡萄糖伴随两个钠离子及 \qty{-60}{mV} 的\omterm{def:b1:membranes-transport:potential}{膜电位}决定；其实现依赖于
\omterm{def:b1:membranes-transport:transporters}{钠钾泵}（基底外侧）与钠–葡萄糖
\omterm{prop:b1:membranes-transport:secondary}{同向转运体}（顶端）。
\end{solution}
